\chapter{Related Work}\label{ch:related_work}

\Cref{ch:introduction} identified a gap between style authoring and tile data preparation - two concerns that remain independent in all existing open-source tooling.
Before developing an optimiser that bridges this gap, this chapter surveys the state of the art along four dimensions that correspond to the levels at which the optimiser will operate.
\cref{sec:rw_tile_data} reviews existing approaches to tile data pruning and schema design.
\cref{sec:rw_tile_encoding} covers tile encoding formats and columnar compression techniques.
\cref{sec:rw_style_optimisations} examines the limited prior work on style-level optimisation.
\cref{sec:rw_benchmarking} surveys rendering performance benchmarking methodologies.

For each area, the discussion identifies what existing tools achieve and their shortcomings, building toward the composite gap that motivates the co-optimisation approach developed in this thesis.
Specifically, tile data pruning motivates the shaving pass, encoding formats motivate the strategy competition, style optimisation motivates the expression and structural passes, and benchmarking motivates the evaluation methodology.

\section{Tile Data Optimisation}\label{sec:rw_tile_data}

The most direct approach to reducing tile transfer cost is to remove data that the rendering style does not reference.

Mapbox developed \texttt{vtshaver}~\cite{mapboxVtshaverTileShaving2018}, a C++/Node.js tool that takes a vector tile and a Mapbox GL style, analyses which layers, features, and properties the style references, and strips everything else.
This reduces tile size and decoding cost, but \texttt{vtshaver} operates exclusively on the data side.
It does not modify the style itself and therefore cannot simplify filter expressions, merge redundant layers, or reduce per-frame rendering overhead.
It also only shaves on a zoom-property-unaware row-collumn approach.

Mapbox also offers \emph{style-optimized vector tiles} as a proprietary server-side feature~\cite{StyleoptimizedVectorTiles}.
When a tile request includes style metadata, the server analyses the style's sources, filters, and zoom ranges, then strips all data not visible for that style.
This is the closest commercial precedent for style-data co-optimisations, but it is opaque, closed-source, and API-only - it cannot be extended, audited, or deployed independently.
The exact set of optimisations is not disclosed and given that MapBox does not have a open research stance (their terms of service explicitly forbid this), reverse-engineering would likely be illegal\footnote{\url{https://media.ccc.de/v/36c3-11089-reflections_on_the_new_reverse_engineering_law}}.

Vilardaga built \texttt{vt-optimizer}~\cite{vilardagaVtoptimizerVectorTile2023}, which drops invisible or unused layers and fields in both the style and data.
The tool also provides an inspection mode that checks conformance to Mapbox's recommended tile size limits (\qty{50}{\kilo\byte} average, \qty{500}{\kilo\byte} maximum).
However, the optimisations are limited to visibility-based pruning.
It does not perform expression-level rewrites, cross-layer merging, or data-driven constant folding.

Tippecanoe~\cite{FeltTippecanoe2026} is the standard open-source tool for building vector tilesets from large GeoJSON or FlatGeobuf~\cite{FlatGeobuf} collections.
It implements several data-level optimisations heuristics: feature dropping based on density budgets (\texttt{--drop-densest-as-needed}), polygon coalescing at low zoom levels, Visvalingam line simplification~\cite{visvalingamLineGeneralisationRepeated1993}, and attribute exclusion.
The line simplification algorithms used throughout the geospatial toolchain - Douglas-Peucker~\cite{douglasALGORITHMSREDUCTIONNUMBER1973} for its computational simplicity and Visvalingam-Whyatt for its area-based perceptual fidelity - are instances of \emph{cartographic generalisation}: the principled reduction of map detail to match the target scale.
Tile data at different zoom levels is, in effect, a discrete multi-scale generalisation hierarchy, though this thesis does not contribute new generalisation algorithms.
These heuristics are effective but purely data-driven: Tippecanoe has no knowledge of which style will render the tiles and therefore cannot tailor its decisions to the rendering workload.

From a cartographic perspective, style-driven tile shaving can be understood as a form of \emph{model generalisation}~\cite{mcmasterGeneralizationDigitalCartography1992}.
The rendering style serves as the generalisation specification, and the pruning advisory selects which data survives the reduction - analogous to how a target scale determines which features are retained in traditional multi-scale generalisation.
The distinction is that classical generalisation selects features by geometric criteria (area, length, density), whereas style-driven shaving selects by attribute criteria (which properties and values the style references).

At a higher level of abstraction, tile schemas themselves embody data-level design decisions.
The OpenMapTiles schema~\cite{VectorTilesAre2025} defines which attributes are included at which zoom levels for OpenStreetMap data.
Wallner et al.~\cite{wallnerOpenSourceVector2022} describe a PostGIS-based pipeline for dynamic vector tile creation with runtime geometry simplification and \ac{MVT} encoding for spatial data infrastructures.
Tile generators such as Planetiler~\cite{OnthegomapPlanetiler2026} and Tilemaker~\cite{TilemakerDIYVector} apply schema-aware rules during the initial tile build step, and Protomaps basemap~\cite{brandonProtomaps} ships a style-specific schema in which column selection is effectively hard-coded into the tile generator.
These systems define the data that enters the pipeline.
The optimisations in this thesis operate downstream, on the data and style that emerge from such pipelines, and do not require re-running the generator.

A complementary framing connects tile shaving directly to relational query processing.
The transformations developed in \cref{ch:methods} that compute a tile pruning advisory from the style and apply it to the tile data are a form of \emph{projection pushdown}~\cite{graefeQueryEvaluationTechniques1993}.
The rendering style is the \enquote{query} (specifying which columns, geometry types, and property values it reads), the tile set is the \enquote{base relation}, and the advisory is the pushdown plan that moves the projection and predicate evaluation upstream of the wire format.
Columnar analytical systems built on Parquet - whose columnar nested-data representation originates in Dremel's record shredding and assembly algorithm~\cite{melnikDremelInteractiveAnalysis} - or ORC rely on precisely this pushdown to avoid reading unreferenced columns.
Abadi et~al.~\cite{abadiMaterializationStrategiesColumnOriented2007} formalised the trade-off between \emph{early} and \emph{late materialisation} in column stores: early materialisation reconstructs full tuples before applying operators, while late materialisation defers reconstruction until the final projection, accessing only the columns that survive filtering.
The thesis's tile shaving is a form of early materialisation (columns are removed at encoding time), while the selective column decoding discussed in future work (\cref{sec:conclusion_future_work}) is a form of late materialisation (columns are skipped at decode time).
Modern query engines extend the pushdown with predicate filters~\cite{abadiDesignImplementationModern2013}.
Modern distributed query engines such as Spark SQL~\cite{armbrustSparkSQLRelational2015} implement predicate and projection pushdown as composable tree-transformation rules applied in batches by a rule-based optimiser - a direct descendant of the extensible rule-based query rewriting architecture introduced by Pirahesh et~al.~\cite{piraheshExtensibleRuleBased1992} in the Starburst system - an architecture structurally parallel to the three-phase pipeline in this thesis.
The thesis's contribution is to lift this DB-internal optimisations across the network boundary between the tile server and the MapLibre renderer, treating the declarative style as a readable query plan.
This framing does not appear in prior map-rendering work - neither \texttt{vtshaver} nor \texttt{vt-optimizer} positions its pruning as query-engine pushdown, and neither tool composes with style-level rewrites to enlarge the pushdown.

\section{Tile Encoding and Format Innovation}\label{sec:rw_tile_encoding}

The dominant wire format for vector tiles is \ac{MVT}~\cite{Vectortilespec21Master}, which uses Protocol Buffers~\cite{ProtocolBuffers} to encode features in a row-oriented layout.
\ac{MVT} has since been adopted as a conformance class in the OGC API-Tiles standard~\cite{OGCAPITiles}, making it the baseline format for any standards-compliant tile delivery.
Its row orientation means that accessing a single property column requires scanning the entire feature table, and its protobuf encoding does not exploit the statistical properties of individual columns.

Tremmel et al.~\cite{tremmelMapLibreTileNext2025} introduced the MapLibre Tile (MLT) format, a columnar alternative that achieves \qtyrange{2}{6}{\times} tile size reduction over MVT.
\ac{MLT} uses column-oriented storage with lightweight compression - delta encoding, run-length encoding, dictionary encoding, and \ac{FSST}~\cite{bonczFSSTFastRandom2020} - and is designed for GPU-accelerated processing.
At the archive layer, Tremmel~\cite{tremmelCOMTILESCASESTUDY2023} proposed COMTiles, a cloud-optimised tile archive format that complements the per-tile encoding of \ac{MLT} by addressing the container and delivery layer for planet-scale tilesets.
However, their work does not consider styles: the encoding strategy is chosen without knowledge of which properties the renderer will actually access, leaving potential for further size reduction through style-aware pruning before encoding.
Removing properties and features changes column cardinalities and value distributions - for instance, pruning rare road classes can convert a high-cardinality string column into one amenable to dictionary or run-length encoding - so the optimal encoding strategy for the shaved tile differs from the original.

The encoding techniques used in \ac{MLT} draw on a rich literature in columnar database compression.
Lemire and Boytsov~\cite{lemireDecodingBillionsIntegers2015} introduced \ac{SIMD}-vectorised integer compression \ac{FastPFOR} achieving billions of integers per second in decompression throughput.
Boncz et al.~\cite{bonczFSSTFastRandom2020} developed \ac{FSST}, a lightweight string compression scheme using a static symbol table that supports random access without decompressing surrounding data.
Kuschewski et al.~\cite{kuschewskiBtrBlocksEfficientColumnar2023} demonstrated with BtrBlocks that cascading multiple encoding schemes (including \ac{FSST}, bit-packing, \ac{RLE}, and dictionary encoding) with automatic per-block strategy selection achieves 2.2$\times$ faster scans than Parquet.
The general finding that no single encoding is universally optimal - and that data-driven strategy selection is necessary - motivates the multi-strategy encoding competition developed in this thesis.

The problem of automatic encoding selection in columnar databases has a substantial lineage.
Abadi et~al.~\cite{abadiIntegratingCompressionExecution2006} introduced decision rules for choosing among RLE, dictionary, and bit-packing based on column properties (cardinality, sort order, run lengths) in the C-Store system.
Damme et~al.~\cite{dammeComprehensiveExperimentalSurvey2019} evaluated dozens of lightweight integer encodings and developed a cost-based selection strategy. Boissier~\cite{ganBaguaScalingDistributed2021} formulated encoding selection as a budget-constrained optimisation problem.
And Zeng et~al.~\cite{zengEmpiricalEvaluationColumnar2023} empirically evaluated Parquet and ORC internals, confirming that dictionary encoding should be a default candidate and that decompression speed matters more than compression ratio for integer encodings.

The thesis's encoding competition differs from these systems by operating at tile granularity (\qtyrange{10}{100}{\kilo\byte}) rather than on large analytical datasets.
We also includes application specific optimsations like spatial sort-order dimension (ID, Hilbert, Morton) that database encoders do not consider.
This competition is an instance of \emph{instance-optimised} system design as proposed by Kraska~\cite{dingSageDBInstanceOptimizedData2022}: rather than committing to a fixed encoding plan, the encoder uses the data distribution of each tile to select optimal physical representations, exactly the approach advocated by the SageDB vision~\cite{dingSageDBInstanceOptimizedData2022} and realised concretely by Ding et~al.~\cite{dingInstanceOptimizedDataLayouts2021}, who showed that instance-optimised data layouts reduce blocks accessed by 93\% on cloud analytics workloads, and at the column-block level by BtrBlocks~\cite{kuschewskiBtrBlocksEfficientColumnar2023}.
Alternative approaches to encoding selection include Fehér et~al.'s~\cite{feherAdaptiveColumnCompression2022} adaptive compression family, which selects encodings at runtime using segment-level measurements, and Cen et~al.'s~\cite{cenLEALearnedEncoding2021} LEA, a learned encoding advisor that uses ML models to predict the best encoding per column - achieving 19\% lower latency and 26\% less space than heuristic selection on TPC-H.
The thesis's trial-based competition is closest in spirit to the profile-guided approach, but substitutes exhaustive evaluation over a bounded candidate space for the learned or heuristic strategies used in database systems.

Gaffuri~\cite{gaffuriWebMappingVector2012} presented one of the earliest academic treatments of vector tile architecture for web maps, proposing a framework integrating tiling, spatial indexing, and generalisation for multi-scale data.
His work established the architectural vocabulary that \ac{MVT} and \ac{MLT} build on, but predates the shift to columnar encoding and does not consider style-awareness as a dimension of the encoding problem - precisely the dimensions that this thesis adds to the tile encoding pipeline.

GeoParquet's columnar layout with Parquet's built-in encoding competition (dictionary, \ac{RLE}, delta, bit-packing per page) is structurally parallel to the \ac{MLT} encoder's per-stream competition.
The key difference is granularity: Parquet optimises for analytical scans over large row groups, whereas \ac{MLT} optimises for tile-granularity access in an interactive rendering pipeline where each tile is a single, small columnar unit (\qtyrange{10}{100}{\kilo\byte}).
Parquet's per-column-chunk metadata overhead would dominate at these tile sizes, which is why a purpose-built tile format remains justified despite the ecosystem support GeoParquet provides.

\section{Style Optimisation}\label{sec:rw_style_optimisations}

In contrast to data-level optimisations, style-level optimisation techniques have attracted comparatively little research attention.

\begin{sloppypar}
The MapLibre style specification includes the utilities \texttt{gl\allowbreak-style\allowbreak-validate} and \texttt{gl\allowbreak-style\allowbreak-migrate}~\cite{MapLibreStyleSpec}, which validate style documents against the specification and mechanically migrate legacy function syntax to the expression language.
These tools are not simplifiers - they are deliberately conservative transformers - but the mechanical migration they produce often contains redundancy that a simplifier could improve.
Deeply nested \texttt{case} chains where a \texttt{match} is sufficent, explicit \texttt{typeof} guards that are trivially true after the migration, and boolean expressions that could be flattened.
A simplification pass is therefore a natural complement to migration, not a replacement for it.
\end{sloppypar}

More broadly, style documents in practice accumulate redundancy through several channels: iterative authoring that leaves superseded properties in place, copy-paste development across layers, automated migration that introduces mechanically correct but verbose expression forms, and multi-author editing where conventions drift over successive contributions.
The result is that production styles routinely contain default-valued properties, unreachable filter branches, and duplicated paint blocks - creating a practical demand for automatic simplification that no existing tool meets.

Stamen Design published \texttt{mapbox-gl-style-remove-defaults}\footnote{\url{https://github.com/stamen/mapbox-gl-style-remove-defaults}}, a tool that removes properties set to their specification default values from Mapbox GL style JSON.
This is the default stripping pass described in \cref{ch:methods}, implemented as a standalone tool.
This ecosystem tooling demonstrates that individual style optimisations are recognised as valuable by practitioners, but no existing tool composes multiple optimisations passes into a pipeline.

Outside the map domain, the most closely related academic work is CSS minification.
Hague et al.~\cite{hagueCSSMinificationConstraint2019} formalised CSS rule merging as a graph optimisations problem: they model rules and properties as a node-weighted bipartite graph with ordering constraints and solve the merging problem using a Max-SAT solver.
Their tool (SatCSS) outperforms six popular CSS minifiers on real-world benchmarks.
The problem structure is directly analogous to the layer merging pass in this thesis: both address merging declarative style rules that differ in selectors (CSS) or filters (map styles) but share properties, subject to ordering constraints imposed by the cascade (CSS) or painter's algorithm (map rendering).

More broadly, rule-based optimisation is a well-studied problem in both databases and compilers.
Pirahesh et~al.~\cite{piraheshExtensibleRuleBased1992} established the paradigm of \emph{rule-based query rewriting} as a distinct optimisation phase in the Starburst system.
A set of declarative transformation rules is applied in priority order to rewrite queries into more efficient forms before cost-based plan selection.
The expression-level passes in this thesis follow the same architecture. Declarative rewrite rules are applied in priority order to style expressions. This makes Starburst the direct database ancestor of the approach.
The modern state of the art for avoiding rule-ordering sensitivity is \emph{equality saturation}~\cite{tateEqualitySaturationNew2010,willseyEggFastExtensible2020}: rather than applying rewrites sequentially and committing to the first match, an e-graph simultaneously represents all equivalent forms and selects the optimal one after saturation.
The rewrite system in this thesis is not confluent and applies rules in a fixed priority order; equality saturation would eliminate this ordering dependence at the cost of significantly higher implementation and runtime complexity.
Since the peephole rules are individually simple and the styles converge within a few fixpoint iterations (\cref{sub:fixpoint} provides the convergence analysis), the simpler sequential approach suffices for the current use case, but equality saturation remains a natural future direction if the rule set grows or ordering sensitivity becomes a practical concern.

\section{Rendering Performance Benchmarking}\label{sec:rw_benchmarking}

Several studies have benchmarked vector map rendering performance.

Netek et al.~\cite{netekPerformanceTestingVector2020} compared vector and raster tile performance across eight scenarios, measuring loading time, data transfer size, and HTTP request count.
They find that vector tiles offer better loading times, though raster tiles may be preferable on very slow connections.

Balla and Gede~\cite{ballaVectorDataRendering2025} benchmarked GeoJSON rendering performance across Leaflet, OpenLayers, MapLibre GL~JS, and Mapbox GL~JS with varying feature counts.
Below 10\,000 features, SVG-based renderers (Leaflet, OpenLayers) are faster; above 50\,000 features, WebGL-based renderers (Mapbox GL~JS, MapLibre GL~JS) dominate due to GPU acceleration.

\begin{sloppypar}
Rechsteiner~\cite{rechsteinerPerformancevergleichOpenSourceVectorTilesServerloesungenZur2024} compared open-source vector tile server implementations.
Since the systems under test offer different feature sets, the evaluation partly compares incomparable aspects and does not address client-side rendering optimisations or style-level improvements.
\end{sloppypar}

Georges et~al.~\cite{georgesStatisticallyRigorousJava2007} and Hoefler and Belli~\cite{hoeflerScientificBenchmarkingParallel2015} established the methodological foundations for statistically rigorous performance evaluation; the specific techniques drawn from their work - bootstrap confidence intervals, warm-up handling, and metric separation - are introduced in \cref{sec:benchmarking}.

These benchmarking studies establish the measurement landscape but focus on comparing \emph{systems} rather than evaluating the impact of \emph{optimisations passes} applied to a fixed system.
This thesis complements them by introducing a cumulative ablation methodology that isolates the contribution of each optimisations pass within a single renderer (MapLibre GL~JS).

\section{Positioning}\label{sec:rw_positioning}

In summary, existing work addresses individual optimisation dimensions in isolation.
Data-level tools such as \texttt{vtshaver}, \texttt{vt-optimizer}, and Tippecanoe prune tile content but leave the style untouched.
Conversely, Stamen's default-stripping tool and the CSS merging work by Hague et al.\ operate purely on style rules without considering the underlying data.
Tremmel and Zink advance tile encoding with the columnar \ac{MLT} format yet select encoding strategies without style knowledge.
The benchmarking studies by Netek et al., Balla and Gede, and Rechsteiner establish measurement baselines but evaluate systems rather than optimisation passes.
Mapbox's proprietary style-optimised tiles are the only precedent that couples style and data, but the approach is closed-source and cannot be inspected, extended, or independently deployed.
The recent formalisation of vector tile delivery in the OGC API-Tiles standard~\cite{OGCAPITiles} underscores the need for optimisation techniques that are compatible with standardised, open infrastructure rather than tied to a single vendor's API.

Critically, no existing open-source tool performs \emph{systematic composition} of style and data optimisations.
Using tile statistics to inform style simplification, and using the simplified style to derive a pruning advisory for the data, so that each level's output feeds forward into the next.
This thesis fills the gap by applying compiler and database optimisation techniques to the map style and tile data as a unified optimisation target, measuring the interaction between levels through a cumulative ablation methodology that isolates each pass's contribution and quantifies whether the composition is additive or synergistic.
The magnitude of this gap is quantified in \cref{ch:experiments}, which demonstrates that composing these techniques reduces load times by up to \qty{75}{\percent} and tile volume by \qty{28}{\percent}.

The following chapter introduces the theoretical foundations that these techniques draw on.
The map rendering architecture that defines the optimisation target, the compiler and query optimisation concepts that inform the rewrite rules, and the encoding strategies that underpin the data-level passes.
